\documentclass[11pt,a4paper]{article}
\usepackage[utf8]{inputenc}
\usepackage[T1]{fontenc}
\usepackage[margin=2.2cm]{geometry}
\usepackage{amsmath,amssymb}
\usepackage{enumitem}
\usepackage{hyperref}

\setlength{\parindent}{0pt}
\setlength{\parskip}{0.75em}

\title{MA087G Mathematical Statistics\\Beginner-Friendly Full Solution Guide\\2025-08-23}
\author{}
\date{}

\begin{document}
\maketitle

\section*{Source notes before solving}
The official exam PDF was image-only. That means normal text extraction found no text, but the rendered page images were readable.

One small source warning:
\begin{itemize}[leftmargin=*]
    \item In Uppgift 3c, the printed expression is visually a little unclear. The sensible reading, also used by student attempts, is \(Y=X_2-2X_1\). I solve that version. If the expression were read literally as \(Y=X_2-2X_2=-X_2\), then \(E(Y)=-2\) still, but \(V(Y)=2\) instead of \(10\).
\end{itemize}

\section*{Uppgift 1}

\subsection*{What this question is really asking}
We have two events, \(A\) and \(B\). An event just means something that may or may not happen.

The question gives probabilities involving \(A^c\) and \(B^c\). The little \(c\) means ``not''. So:
\[
A^c=\text{not } A,\qquad B^c=\text{not } B.
\]

We must find \(P(A)\) and \(P(B)\), meaning the probability that \(A\) happens and the probability that \(B\) happens.

\subsection*{How to recognize this type of question}
The clues are:
\begin{itemize}[leftmargin=*]
    \item The question contains events \(A\), \(B\), \(A^c\), and \(B^c\).
    \item It gives conditional probabilities like \(P(A^c\mid B^c)\).
    \item It gives a union probability \(P(A^c\cup B^c)\).
\end{itemize}

So this is a probability algebra question. We should translate each line into formulas and solve for the missing probabilities.

\subsection*{Tiny theory needed}
The symbol \(P(A^c\mid B^c)\) means:

\[
\text{probability that \(A\) does not happen, given that \(B\) does not happen.}
\]

The formula for conditional probability is:
\[
P(A^c\mid B^c)=\frac{P(A^c\cap B^c)}{P(B^c)}.
\]

The symbol \(\cap\) means ``and''. So \(A^c\cap B^c\) means ``not \(A\) and not \(B\)''.

The symbol \(\cup\) means ``or''. The formula for an ``or'' probability is:
\[
P(A^c\cup B^c)=P(A^c)+P(B^c)-P(A^c\cap B^c).
\]

We subtract the overlap once because otherwise it gets counted twice.

\subsection*{Formula and what everything means}
We will use:
\[
P(A^c\mid B^c)=\frac{P(A^c\cap B^c)}{P(B^c)}
\]
and
\[
P(B^c\mid A^c)=\frac{P(A^c\cap B^c)}{P(A^c)}.
\]

Here:
\begin{itemize}[leftmargin=*]
    \item \(P(A^c\cap B^c)\) is the probability that neither \(A\) nor \(B\) happens.
    \item \(P(A^c)\) is the probability that \(A\) does not happen.
    \item \(P(B^c)\) is the probability that \(B\) does not happen.
    \item \(P(A)=1-P(A^c)\), because \(A\) either happens or does not happen.
    \item \(P(B)=1-P(B^c)\), for the same reason.
\end{itemize}

\subsection*{Step-by-step solution}
Let
\[
x=P(A^c\cap B^c).
\]
This is just a shortcut so we do not have to write the long expression all the time.

The first given probability is:
\[
P(A^c\mid B^c)=0.6.
\]
Use the conditional probability formula:
\begin{align*}
0.6
&=\frac{P(A^c\cap B^c)}{P(B^c)}\\
&=\frac{x}{P(B^c)}.
\end{align*}
This says \(x\) is \(60\%\) of \(P(B^c)\). Solve for \(P(B^c)\):
\begin{align*}
0.6P(B^c)&=x\\
P(B^c)&=\frac{x}{0.6}
=\frac{5x}{3}.
\end{align*}

The second given probability is:
\[
P(B^c\mid A^c)=0.4.
\]
Again use conditional probability:
\begin{align*}
0.4
&=\frac{P(A^c\cap B^c)}{P(A^c)}\\
&=\frac{x}{P(A^c)}.
\end{align*}
Now solve for \(P(A^c)\):
\begin{align*}
0.4P(A^c)&=x\\
P(A^c)&=\frac{x}{0.4}
=\frac{5x}{2}.
\end{align*}

The third given probability is:
\[
P(A^c\cup B^c)=0.8.
\]
Now use the union formula:
\begin{align*}
P(A^c\cup B^c)
&=P(A^c)+P(B^c)-P(A^c\cap B^c).
\end{align*}
Insert what we know:
\begin{align*}
0.8
&=\frac{5x}{2}+\frac{5x}{3}-x.
\end{align*}
Now put everything over the same denominator, \(6\):
\begin{align*}
0.8
&=\frac{15x}{6}+\frac{10x}{6}-\frac{6x}{6}\\
&=\frac{19x}{6}.
\end{align*}
Solve for \(x\):
\begin{align*}
x
&=0.8\cdot \frac{6}{19}\\
&=\frac{4}{5}\cdot \frac{6}{19}\\
&=\frac{24}{95}.
\end{align*}

Now find \(P(A^c)\):
\begin{align*}
P(A^c)
&=\frac{5x}{2}\\
&=\frac{5}{2}\cdot \frac{24}{95}\\
&=\frac{60}{95}
=\frac{12}{19}.
\end{align*}
So
\begin{align*}
P(A)
&=1-P(A^c)\\
&=1-\frac{12}{19}\\
&=\frac{7}{19}.
\end{align*}

Now find \(P(B^c)\):
\begin{align*}
P(B^c)
&=\frac{5x}{3}\\
&=\frac{5}{3}\cdot \frac{24}{95}\\
&=\frac{40}{95}
=\frac{8}{19}.
\end{align*}
So
\begin{align*}
P(B)
&=1-P(B^c)\\
&=1-\frac{8}{19}\\
&=\frac{11}{19}.
\end{align*}

\subsection*{Answer}
\[
\boxed{P(A)=\frac{7}{19}\approx 0.368}
\]
\[
\boxed{P(B)=\frac{11}{19}\approx 0.579}
\]

\subsection*{What to memorize for the exam}
When you see conditional probabilities with complements, introduce a small unknown for the overlap, for example:
\[
x=P(A^c\cap B^c).
\]
Then rewrite the conditional probabilities using fractions.

\subsection*{Common trap}
Do not think \(P(A^c\mid B^c)\) means \(P(A^c)\cdot P(B^c)\). Conditional probability is a fraction, not multiplication.

\subsection*{Repeated exam pattern}
This is the same basic pattern as 2026-01, Uppgift 1: conditional probability with complements. The numbers and wording changed, but the method is the same: translate the probability statements, use conditional probability formulas, and solve for the missing probabilities.

\section*{Uppgift 2}

\subsection*{What this question is really asking}
The question gives a function \(f(x)\). We must make it into a probability density function.

A density function is used for a continuous random variable, meaning a variable that can take values on an interval, not just separate values like \(0,1,2\).

We need to:
\begin{itemize}[leftmargin=*]
    \item find \(a\), the missing constant,
    \item find the distribution function \(F_X(x)\),
    \item calculate \(P(X>1)\),
    \item calculate expected value and standard deviation.
\end{itemize}

\subsection*{How to recognize this type of question}
The clues are:
\begin{itemize}[leftmargin=*]
    \item The question says ``täthetsfunktion'', which means density function.
    \item The function is given on an interval \(0<x<3\).
    \item It asks for \(F_X\), which is the distribution function/CDF.
    \item It asks for expected value and standard deviation.
\end{itemize}

So we should use integrals. For continuous variables, sums become integrals.

\subsection*{Tiny theory needed}
The density is:
\[
f(x)=
\begin{cases}
ax^2, & 0<x<3,\\
0, & \text{otherwise}.
\end{cases}
\]

For \(f\) to be a density, the total area under the curve must be \(1\). That is because total probability is always \(1\).

So we need:
\[
\int_{-\infty}^{\infty} f(x)\,dx=1.
\]
But \(f(x)=0\) outside \(0<x<3\), so we only integrate from \(0\) to \(3\).

\subsection*{Formula and what everything means}
Density condition:
\[
\int_0^3 f(x)\,dx=1.
\]
Here:
\begin{itemize}[leftmargin=*]
    \item \(f(x)\) is the density.
    \item The integral means area under the curve.
    \item The answer must be \(1\), because all probabilities together add to \(1\).
\end{itemize}

Distribution function:
\[
F_X(x)=P(X\le x)=\int_{-\infty}^{x} f(t)\,dt.
\]
Here:
\begin{itemize}[leftmargin=*]
    \item \(F_X(x)\) means probability that \(X\) is at most \(x\).
    \item \(t\) is just a dummy integration letter. It plays the same role as \(x\), but avoids confusion.
\end{itemize}

Expected value:
\[
E(X)=\int x f(x)\,dx.
\]
This is the average value of \(X\) in the long run.

Variance:
\[
V(X)=E(X^2)-[E(X)]^2.
\]
Standard deviation:
\[
D(X)=\sqrt{V(X)}.
\]
The standard deviation tells us roughly how spread out \(X\) is.

\subsection*{Step-by-step solution}
\textbf{Part a: Find \(a\).}

Use the density condition:
\begin{align*}
\int_0^3 ax^2\,dx&=1.
\end{align*}
The constant \(a\) can be moved outside the integral:
\begin{align*}
a\int_0^3 x^2\,dx&=1.
\end{align*}
Now integrate \(x^2\):
\begin{align*}
\int x^2\,dx=\frac{x^3}{3}.
\end{align*}
So:
\begin{align*}
a\left[\frac{x^3}{3}\right]_0^3&=1\\
a\left(\frac{3^3}{3}-\frac{0^3}{3}\right)&=1\\
a\left(\frac{27}{3}\right)&=1\\
9a&=1\\
a&=\frac{1}{9}.
\end{align*}

\textbf{Part b: Find \(F_X(x)\).}

There are three regions:
\begin{itemize}[leftmargin=*]
    \item If \(x\le 0\), then we are before the density starts, so \(F_X(x)=0\).
    \item If \(0<x<3\), integrate from \(0\) to \(x\).
    \item If \(x\ge 3\), we have passed the whole density, so \(F_X(x)=1\).
\end{itemize}

For \(0<x<3\):
\begin{align*}
F_X(x)
&=\int_0^x \frac{1}{9}t^2\,dt\\
&=\frac{1}{9}\left[\frac{t^3}{3}\right]_0^x\\
&=\frac{1}{9}\left(\frac{x^3}{3}-0\right)\\
&=\frac{x^3}{27}.
\end{align*}

So:
\[
F_X(x)=
\begin{cases}
0, & x\le 0,\\[1mm]
\dfrac{x^3}{27}, & 0<x<3,\\[2mm]
1, & x\ge 3.
\end{cases}
\]

\textbf{Part c: Calculate \(P(X>1)\).}

The CDF gives \(P(X\le 1)\), so use:
\[
P(X>1)=1-P(X\le 1)=1-F_X(1).
\]
Now insert \(x=1\):
\begin{align*}
F_X(1)&=\frac{1^3}{27}=\frac{1}{27}.
\end{align*}
Therefore:
\begin{align*}
P(X>1)
&=1-\frac{1}{27}\\
&=\frac{27}{27}-\frac{1}{27}\\
&=\frac{26}{27}.
\end{align*}

\textbf{Part d: Expected value and standard deviation.}

First calculate \(E(X)\):
\begin{align*}
E(X)
&=\int_0^3 x f(x)\,dx\\
&=\int_0^3 x\cdot \frac{1}{9}x^2\,dx\\
&=\frac{1}{9}\int_0^3 x^3\,dx\\
&=\frac{1}{9}\left[\frac{x^4}{4}\right]_0^3\\
&=\frac{1}{9}\cdot \frac{3^4}{4}\\
&=\frac{1}{9}\cdot \frac{81}{4}\\
&=\frac{9}{4}.
\end{align*}
So the average value is \(9/4=2.25\).

Now calculate \(E(X^2)\), because we need it for the variance:
\begin{align*}
E(X^2)
&=\int_0^3 x^2 f(x)\,dx\\
&=\int_0^3 x^2\cdot \frac{1}{9}x^2\,dx\\
&=\frac{1}{9}\int_0^3 x^4\,dx\\
&=\frac{1}{9}\left[\frac{x^5}{5}\right]_0^3\\
&=\frac{1}{9}\cdot \frac{3^5}{5}\\
&=\frac{1}{9}\cdot \frac{243}{5}\\
&=\frac{27}{5}.
\end{align*}

Now variance:
\begin{align*}
V(X)
&=E(X^2)-[E(X)]^2\\
&=\frac{27}{5}-\left(\frac{9}{4}\right)^2\\
&=\frac{27}{5}-\frac{81}{16}.
\end{align*}
Use common denominator \(80\):
\begin{align*}
V(X)
&=\frac{432}{80}-\frac{405}{80}\\
&=\frac{27}{80}.
\end{align*}
Standard deviation:
\begin{align*}
D(X)
&=\sqrt{V(X)}\\
&=\sqrt{\frac{27}{80}}\\
&\approx 0.581.
\end{align*}

\subsection*{Answer}
\[
\boxed{a=\frac{1}{9}}
\]
\[
\boxed{
F_X(x)=
\begin{cases}
0, & x\le 0,\\
\frac{x^3}{27}, & 0<x<3,\\
1, & x\ge 3
\end{cases}}
\]
\[
\boxed{P(X>1)=\frac{26}{27}\approx 0.963}
\]
\[
\boxed{E(X)=\frac{9}{4}=2.25},\qquad
\boxed{D(X)=\sqrt{\frac{27}{80}}\approx 0.581.}
\]

\subsection*{What to memorize for the exam}
For density questions:
\begin{itemize}[leftmargin=*]
    \item First make the total area equal to \(1\).
    \item Then build the CDF by integrating the density.
    \item For \(P(X>c)\), use \(1-F(c)\).
    \item For standard deviation, calculate variance first.
\end{itemize}

\subsection*{Common trap}
Do not forget the CDF pieces outside the interval. A CDF is not only \(\frac{x^3}{27}\). It must also be \(0\) before the interval and \(1\) after the interval.

\subsection*{Repeated exam pattern}
This is the same basic pattern as 2026-01, Uppgift 2: continuous density normalization, CDF, probability, expected value, and standard deviation. The density function changed, but the workflow is the same.

\section*{Uppgift 3}

\subsection*{What this question is really asking}
We have two independent Poisson random variables:
\[
X_1,\;X_2\sim \operatorname{Po}(2).
\]
This means both variables count something, and the average count is \(2\).

We need to calculate:
\begin{itemize}[leftmargin=*]
    \item one exact Poisson probability,
    \item one interval Poisson probability,
    \item a probability involving the maximum of \(X_1\) and \(X_2\),
    \item expected value and variance for a linear combination.
\end{itemize}

\subsection*{How to recognize this type of question}
The clues are:
\begin{itemize}[leftmargin=*]
    \item The notation \(\operatorname{Po}(\mu)\) means Poisson distribution.
    \item The question says \(\mu=2\), so the Poisson mean is \(2\).
    \item It says \(X_1\) and \(X_2\) are independent, so probabilities involving both variables can be multiplied.
    \item The word \(\max\) means the largest of the two values.
\end{itemize}

\subsection*{Tiny theory needed}
For a Poisson variable \(X\sim\operatorname{Po}(\mu)\), the probability of exactly \(k\) is:
\[
P(X=k)=\frac{\mu^k}{k!}e^{-\mu}.
\]

For Poisson:
\[
E(X)=\mu,\qquad V(X)=\mu.
\]
So for \(\mu=2\):
\[
E(X_1)=E(X_2)=2,\qquad V(X_1)=V(X_2)=2.
\]

If \(X_1\) and \(X_2\) are independent, then:
\[
P(X_1\le 2,\;X_2\le 2)=P(X_1\le 2)P(X_2\le 2).
\]

\subsection*{Formula and what everything means}
Poisson formula:
\[
P(X=k)=\frac{\mu^k}{k!}e^{-\mu}.
\]
Here:
\begin{itemize}[leftmargin=*]
    \item \(k\) is the exact number we want.
    \item \(\mu\) is the average count.
    \item \(k!\) means \(k\cdot(k-1)\cdots 1\). For example, \(3!=3\cdot2\cdot1=6\).
    \item \(e\) is the calculator number \(2.71828...\).
\end{itemize}

For variance of a linear combination:
\[
V(aX+bY)=a^2V(X)+b^2V(Y)
\]
when \(X\) and \(Y\) are independent. The coefficients are squared, so negative signs disappear in variance.

\subsection*{Step-by-step solution}
\textbf{Part a: Calculate \(P(X_1=1)\).}

Here \(k=1\) and \(\mu=2\):
\begin{align*}
P(X_1=1)
&=\frac{2^1}{1!}e^{-2}\\
&=2e^{-2}\\
&\approx 0.271.
\end{align*}

Now calculate \(P(1\le X_1<4)\). Since \(X_1\) is a count, this means:
\[
X_1=1,\;2,\;\text{or }3.
\]
So:
\begin{align*}
P(1\le X_1<4)
&=P(X_1=1)+P(X_1=2)+P(X_1=3).
\end{align*}
Calculate each part:
\begin{align*}
P(X_1=1)&=\frac{2^1}{1!}e^{-2}=2e^{-2},\\
P(X_1=2)&=\frac{2^2}{2!}e^{-2}=\frac{4}{2}e^{-2}=2e^{-2},\\
P(X_1=3)&=\frac{2^3}{3!}e^{-2}=\frac{8}{6}e^{-2}=\frac{4}{3}e^{-2}.
\end{align*}
Add them:
\begin{align*}
P(1\le X_1<4)
&=\left(2+2+\frac{4}{3}\right)e^{-2}\\
&=\frac{16}{3}e^{-2}\\
&\approx 0.722.
\end{align*}

\textbf{Part b: Calculate \(P(\max(X_1,X_2)\le 2)\).}

The maximum is at most \(2\) only if both variables are at most \(2\). So:
\[
\max(X_1,X_2)\le 2
\]
means:
\[
X_1\le 2 \text{ and } X_2\le 2.
\]
Because \(X_1\) and \(X_2\) are independent:
\[
P(\max(X_1,X_2)\le 2)=P(X_1\le 2)P(X_2\le 2).
\]
They have the same distribution, so these two probabilities are equal:
\[
P(\max(X_1,X_2)\le 2)=[P(X_1\le 2)]^2.
\]

Now calculate \(P(X_1\le 2)\):
\begin{align*}
P(X_1\le 2)
&=P(X_1=0)+P(X_1=1)+P(X_1=2)\\
&=\frac{2^0}{0!}e^{-2}+\frac{2^1}{1!}e^{-2}+\frac{2^2}{2!}e^{-2}\\
&=(1+2+2)e^{-2}\\
&=5e^{-2}.
\end{align*}
Therefore:
\begin{align*}
P(\max(X_1,X_2)\le 2)
&=(5e^{-2})^2\\
&=25e^{-4}\\
&\approx 0.458.
\end{align*}

\textbf{Part c: Calculate \(E(Y)\) and \(V(Y)\).}

Using the intended reading \(Y=X_2-2X_1\):
\[
Y=X_2-2X_1.
\]

Expected value is linear. That means we can move constants outside:
\begin{align*}
E(Y)
&=E(X_2-2X_1)\\
&=E(X_2)-2E(X_1).
\end{align*}
Since both means are \(2\):
\begin{align*}
E(Y)
&=2-2\cdot 2\\
&=2-4\\
&=-2.
\end{align*}

For variance, we square the coefficients:
\begin{align*}
V(Y)
&=V(X_2-2X_1)\\
&=V(X_2)+(-2)^2V(X_1).
\end{align*}
The minus sign disappears because \((-2)^2=4\):
\begin{align*}
V(Y)
&=2+4\cdot 2\\
&=2+8\\
&=10.
\end{align*}

\subsection*{Answer}
\[
\boxed{P(X_1=1)=2e^{-2}\approx 0.271}
\]
\[
\boxed{P(1\le X_1<4)=\frac{16}{3}e^{-2}\approx 0.722}
\]
\[
\boxed{P(\max(X_1,X_2)\le 2)=25e^{-4}\approx 0.458}
\]
\[
\boxed{E(Y)=-2,\qquad V(Y)=10}
\]
for the intended reading \(Y=X_2-2X_1\).

\subsection*{What to memorize for the exam}
For Poisson:
\[
P(X=k)=\frac{\mu^k}{k!}e^{-\mu},\qquad E(X)=\mu,\qquad V(X)=\mu.
\]
For a maximum:
\[
\max(X_1,X_2)\le c
\]
means both variables are \(\le c\).

\subsection*{Common trap}
Variance can never be negative. When calculating \(V(X_2-2X_1)\), the \(-2\) must be squared:
\[
(-2)^2=4.
\]
So the variance contribution is \(4V(X_1)\), not \(-2V(X_1)\).

\subsection*{Repeated exam pattern}
The Poisson parts are a new pattern in the processed folders so far. The linear-combination variance idea is related to 2026-01, Uppgift 4b: coefficients must be squared when working with variance.

\section*{Uppgift 4}

\subsection*{What this question is really asking}
This is a maximum likelihood question, usually shortened to ML.

We are given a density:
\[
f_X(x)=\theta x^{\theta-1},\qquad 0<x<1.
\]
The unknown parameter is \(\theta\). The question wants the ML estimate of \(\theta\), meaning the value of \(\theta\) that makes the observed sample most likely.

\subsection*{How to recognize this type of question}
The clue is the phrase:
\[
\text{``Ange ML-skattningen av \(\theta\)''}.
\]
That tells us immediately:
\begin{itemize}[leftmargin=*]
    \item write the likelihood,
    \item take the logarithm,
    \item differentiate,
    \item set the derivative equal to zero,
    \item solve for \(\theta\).
\end{itemize}

\subsection*{Tiny theory needed}
The likelihood is the product of the density values at all observed data points:
\[
L(\theta)=\prod_{i=1}^n f_X(x_i).
\]
The product sign \(\prod\) just means multiply:
\[
\prod_{i=1}^n f_X(x_i)=f_X(x_1)f_X(x_2)\cdots f_X(x_n).
\]

We take logarithms because products are annoying, but logs turn products into sums. This makes differentiation much easier.

\subsection*{Formula and what everything means}
The density is:
\[
f_X(x_i)=\theta x_i^{\theta-1}.
\]
Here:
\begin{itemize}[leftmargin=*]
    \item \(x_i\) is one observed value in the sample.
    \item \(n\) is the number of observations.
    \item \(\theta\) is the unknown number we are estimating.
    \item \(L(\theta)\) is the likelihood, a function of \(\theta\).
\end{itemize}

ML recipe:
\[
\ell(\theta)=\ln L(\theta),
\]
then solve:
\[
\ell'(\theta)=0.
\]

\subsection*{Step-by-step solution}
Start with the likelihood:
\begin{align*}
L(\theta)
&=\prod_{i=1}^n \theta x_i^{\theta-1}.
\end{align*}
There is one \(\theta\) for every observation, so multiplying them gives \(\theta^n\):
\begin{align*}
L(\theta)
&=\theta^n\prod_{i=1}^n x_i^{\theta-1}.
\end{align*}

Now take the natural log:
\begin{align*}
\ell(\theta)
&=\ln L(\theta)\\
&=\ln\left(\theta^n\prod_{i=1}^n x_i^{\theta-1}\right).
\end{align*}
Use log rules:
\begin{align*}
\ell(\theta)
&=n\ln\theta+\sum_{i=1}^n \ln\left(x_i^{\theta-1}\right).
\end{align*}
Another log rule says \(\ln(a^b)=b\ln(a)\):
\begin{align*}
\ell(\theta)
&=n\ln\theta+\sum_{i=1}^n(\theta-1)\ln x_i\\
&=n\ln\theta+(\theta-1)\sum_{i=1}^n \ln x_i.
\end{align*}

Now differentiate with respect to \(\theta\):
\begin{align*}
\ell'(\theta)
&=\frac{n}{\theta}+\sum_{i=1}^n \ln x_i.
\end{align*}
Set the derivative equal to zero:
\begin{align*}
\frac{n}{\theta}+\sum_{i=1}^n \ln x_i&=0.
\end{align*}
Move the sum to the other side:
\begin{align*}
\frac{n}{\theta}&=-\sum_{i=1}^n \ln x_i.
\end{align*}
Now solve for \(\theta\):
\begin{align*}
\theta
&=\frac{n}{-\sum_{i=1}^n \ln x_i}\\
&=-\frac{n}{\sum_{i=1}^n \ln x_i}.
\end{align*}

Since every \(x_i\) is between \(0\) and \(1\), every \(\ln x_i\) is negative. So the denominator is negative, and the minus sign makes the estimate positive.

\subsection*{Answer}
\[
\boxed{\hat{\theta}_{ML}=-\frac{n}{\sum_{i=1}^n \ln x_i}}
\]

\subsection*{What to memorize for the exam}
For ML questions:
\[
\text{likelihood} \rightarrow \text{log-likelihood} \rightarrow \text{differentiate} \rightarrow \text{set equal to zero}.
\]
That is the exam routine almost every time.

\subsection*{Common trap}
Do not stop after writing the likelihood. Most points come from taking the log, differentiating, and solving for the parameter.

\subsection*{Repeated exam pattern}
This is the same broad ML pattern as 2026-01, Uppgift 5. The distribution is different, but the solving recipe is the same: build \(L\), take \(\ln L\), differentiate, solve.

\section*{Uppgift 5}

\subsection*{What this question is really asking}
We have 8 measurements that are assumed to come from a normal distribution:
\[
N(\mu,\sigma).
\]
The question asks about \(\sigma^2\) and \(\sigma\).

Here:
\begin{itemize}[leftmargin=*]
    \item \(\sigma^2\) is the variance, which measures spread.
    \item \(\sigma\) is the standard deviation, which is the square root of the variance.
\end{itemize}

Part a asks for point estimates, meaning one best guess number for each.

Part b asks for a 90\% confidence interval for \(\sigma\), meaning an interval of plausible values for the true standard deviation.

\subsection*{How to recognize this type of question}
The clues are:
\begin{itemize}[leftmargin=*]
    \item The data are assumed normal.
    \item The question asks about \(\sigma\), not \(\mu\).
    \item A confidence interval for \(\sigma\) uses the chi-square distribution, not the normal distribution.
\end{itemize}

\subsection*{Tiny theory needed}
For sample variance, we use:
\[
s^2=\frac{1}{n-1}\left(\sum_{i=1}^n x_i^2-n\bar{x}^2\right).
\]

This formula uses:
\begin{itemize}[leftmargin=*]
    \item \(\sum x_i\), the sum of the observations,
    \item \(\sum x_i^2\), the sum of squared observations,
    \item \(\bar{x}\), the sample mean,
    \item \(n-1\), because this is the usual unbiased estimate of variance.
\end{itemize}

For a confidence interval for \(\sigma\), we first make an interval for \(\sigma^2\), then take square roots.

\subsection*{Formula and what everything means}
Sample mean:
\[
\bar{x}=\frac{\sum x_i}{n}.
\]
Here \(\bar{x}\) is the average of the observed data.

Sample variance:
\[
s^2=\frac{\sum x_i^2-n\bar{x}^2}{n-1}.
\]
Here \(s^2\) estimates \(\sigma^2\).

Confidence interval for \(\sigma^2\):
\[
\left(
\frac{(n-1)s^2}{\chi^2_{0.95,n-1}},
\frac{(n-1)s^2}{\chi^2_{0.05,n-1}}
\right).
\]
For a 90\% interval, we put \(5\%\) in each tail, so we use \(0.95\) and \(0.05\) chi-square values.

\subsection*{Step-by-step solution}
We are given:
\[
n=8,\qquad \sum_{i=1}^8 x_i=114.1,\qquad \sum_{i=1}^8 x_i^2=1629.45.
\]

\textbf{Part a: Point estimate of \(\sigma^2\) and \(\sigma\).}

First calculate the sample mean:
\begin{align*}
\bar{x}
&=\frac{114.1}{8}\\
&=14.2625.
\end{align*}

Now calculate the numerator of the variance:
\begin{align*}
\sum x_i^2-n\bar{x}^2
&=1629.45-8(14.2625)^2.
\end{align*}
Calculate the square:
\[
(14.2625)^2=203.41890625.
\]
Multiply by \(8\):
\[
8(14.2625)^2=1627.35125.
\]
Subtract:
\begin{align*}
1629.45-1627.35125&=2.09875.
\end{align*}
Now divide by \(n-1=7\):
\begin{align*}
s^2
&=\frac{2.09875}{7}\\
&\approx 0.2998.
\end{align*}
So:
\[
\hat{\sigma}^2=s^2\approx 0.2998.
\]
Now take the square root to estimate \(\sigma\):
\begin{align*}
s
&=\sqrt{0.2998}\\
&\approx 0.5476.
\end{align*}

\textbf{Part b: 90\% confidence interval for \(\sigma\).}

We use \(n-1=7\) degrees of freedom.

From a chi-square table:
\[
\chi^2_{0.95,7}\approx 14.067,\qquad
\chi^2_{0.05,7}\approx 2.167.
\]

Also:
\[
(n-1)s^2=7\cdot 0.2998=2.09875.
\]
So the interval for \(\sigma^2\) is:
\begin{align*}
\left(
\frac{2.09875}{14.067},
\frac{2.09875}{2.167}
\right)
&\approx (0.1492,\;0.9683).
\end{align*}

But the question asks for \(\sigma\), not \(\sigma^2\). So take square roots of both endpoints:
\begin{align*}
\sigma
&\in \left(\sqrt{0.1492},\sqrt{0.9683}\right)\\
&\approx (0.386,\;0.984).
\end{align*}

\subsection*{Answer}
\[
\boxed{\hat{\sigma}^2=s^2\approx 0.2998}
\]
\[
\boxed{\hat{\sigma}=s\approx 0.5476}
\]
\[
\boxed{\text{90\% confidence interval for }\sigma:\ (0.386,\;0.984)}
\]

\subsection*{What to memorize for the exam}
If the question asks for an interval for \(\sigma\) or \(\sigma^2\) from normal data, think:
\[
\text{chi-square distribution}.
\]
Do not use a \(z\)-interval or \(t\)-interval for \(\sigma\).

\subsection*{Common trap}
A very common mistake is to calculate a confidence interval for the mean \(\mu\). But this question asks for \(\sigma\), the spread. That is a different formula.

\subsection*{Repeated exam pattern}
No exact earlier matching pattern found yet. Track this as a new pattern: point estimate and chi-square confidence interval for normal standard deviation.

\section*{Uppgift 6}

\subsection*{What this question is really asking}
This question is about testing whether the average passenger weight has increased.

Old belief:
\[
\mu=80.
\]
Now we want to test whether:
\[
\mu>80.
\]

We know the population standard deviation:
\[
\sigma=10.
\]
Because \(\sigma\) is known, this is a \(z\)-test.

\subsection*{How to recognize this type of question}
The clues are:
\begin{itemize}[leftmargin=*]
    \item It says ``testa'', so this is a hypothesis test.
    \item It asks whether the mean has increased, so the alternative is \(>\).
    \item It gives \(\sigma=10\), so use a normal \(z\)-test instead of a \(t\)-test.
    \item Part c asks how large \(n\) must be to detect an increase with 80\% probability. That is a power/sample-size question.
\end{itemize}

\subsection*{Tiny theory needed}
A hypothesis test compares two ideas:
\begin{itemize}[leftmargin=*]
    \item \(H_0\): the default idea, usually ``no increase''.
    \item \(H_1\): the thing we are looking for evidence of.
\end{itemize}

Here we are looking for evidence that the mean is larger than \(80\), so:
\[
H_0:\mu=80,\qquad H_1:\mu>80.
\]

The test statistic tells us how many standard errors the observed mean is above \(80\).

\subsection*{Formula and what everything means}
The \(z\)-test statistic is:
\[
z=\frac{\bar{x}-\mu_0}{\sigma/\sqrt{n}}.
\]
Here:
\begin{itemize}[leftmargin=*]
    \item \(\bar{x}\) is the sample average.
    \item \(\mu_0\) is the mean in the null hypothesis, here \(80\).
    \item \(\sigma\) is the known standard deviation, here \(10\).
    \item \(n\) is the sample size.
    \item \(\sigma/\sqrt{n}\) is the standard error, meaning the typical spread of the sample average.
\end{itemize}

For a one-sided 5\% test:
\[
z_{0.95}\approx 1.645.
\]
We reject \(H_0\) if the test statistic is bigger than \(1.645\).

\subsection*{Step-by-step solution}
\textbf{Part a: Set up the hypotheses.}

The old mean is \(80\), and we are checking if the mean has increased. So:
\[
H_0:\mu=80
\]
and
\[
H_1:\mu>80.
\]

The significance level is \(5\%\). That means we allow a \(5\%\) risk of falsely saying there is an increase when there is not.

\textbf{Part b: Perform the test for \(n=100\) and \(\bar{x}=83\).}

Insert the values:
\[
\bar{x}=83,\qquad \mu_0=80,\qquad \sigma=10,\qquad n=100.
\]

First calculate the standard error:
\begin{align*}
\frac{\sigma}{\sqrt{n}}
&=\frac{10}{\sqrt{100}}\\
&=\frac{10}{10}\\
&=1.
\end{align*}

Now calculate the \(z\)-value:
\begin{align*}
z
&=\frac{\bar{x}-\mu_0}{\sigma/\sqrt{n}}\\
&=\frac{83-80}{1}\\
&=3.
\end{align*}

Compare with the critical value:
\[
1.645.
\]
Since:
\[
3>1.645,
\]
we reject \(H_0\).

In plain language: the sample mean \(83\) is far enough above \(80\) that we say there is evidence the average weight has increased.

\textbf{Part c: Find \(n\) for 80\% probability of detecting \(\mu=82\).}

The test rejects when:
\[
\bar{X}>80+1.645\frac{10}{\sqrt{n}}.
\]
This is the rejection boundary. The sample mean must be above this number to reject \(H_0\).

Now assume the true mean is actually:
\[
\mu=82.
\]
Then:
\[
\bar{X}\sim N\left(82,\frac{10}{\sqrt{n}}\right).
\]
We want:
\[
P(\text{reject }H_0)>0.80.
\]
Using the rejection rule:
\[
P\left(\bar{X}>80+1.645\frac{10}{\sqrt{n}}\right)>0.80.
\]

Standardize this probability. Subtract the true mean \(82\), then divide by the standard error \(10/\sqrt{n}\):
\[
P\left(
Z>
\frac{80+1.645\frac{10}{\sqrt{n}}-82}{10/\sqrt{n}}
\right)>0.80.
\]

The problem tells us:
\[
\Phi(0.84)\approx 0.80.
\]
That means:
\[
P(Z>-0.84)\approx 0.80.
\]
So we need the cutoff inside the probability to be at most \(-0.84\):
\[
\frac{80+1.645\frac{10}{\sqrt{n}}-82}{10/\sqrt{n}}\le -0.84.
\]
Simplify the top:
\[
80-82=-2.
\]
So:
\[
\frac{-2+1.645\frac{10}{\sqrt{n}}}{10/\sqrt{n}}\le -0.84.
\]
Split the fraction into two parts:
\begin{align*}
\frac{-2}{10/\sqrt{n}}+1.645&\le -0.84.
\end{align*}
Since dividing by \(10/\sqrt{n}\) is the same as multiplying by \(\sqrt{n}/10\):
\begin{align*}
-\frac{2\sqrt{n}}{10}+1.645&\le -0.84\\
-0.2\sqrt{n}+1.645&\le -0.84.
\end{align*}
Move \(1.645\) to the right:
\begin{align*}
-0.2\sqrt{n}&\le -0.84-1.645\\
-0.2\sqrt{n}&\le -2.485.
\end{align*}
Divide by \(-0.2\). Because we divide by a negative number, the inequality sign flips:
\begin{align*}
\sqrt{n}&\ge 12.425.
\end{align*}
Square both sides:
\begin{align*}
n&\ge (12.425)^2\\
n&\ge 154.38.
\end{align*}
The sample size must be a whole number, and it must be large enough, so round up:
\[
n=155.
\]

\subsection*{Answer}
\[
\boxed{H_0:\mu=80,\qquad H_1:\mu>80}
\]
\[
\boxed{z=3}
\]
\[
\boxed{\text{Reject }H_0\text{ at the 5\% level. There is evidence that the mean has increased.}}
\]
\[
\boxed{n=155\text{ is needed for about 80\% power when }\mu=82.}
\]

\subsection*{What to memorize for the exam}
If the population standard deviation \(\sigma\) is known, use:
\[
z=\frac{\bar{x}-\mu_0}{\sigma/\sqrt{n}}.
\]
If the question says ``has increased'', use a right-tailed test:
\[
H_1:\mu>\mu_0.
\]

\subsection*{Common trap}
In part c, do not plug in \(\bar{x}=83\). Part c is not about the old sample anymore. It asks how large \(n\) must be if the true mean is \(82\).

\subsection*{Repeated exam pattern}
This is related to 2026-01, Uppgift 6 because both are hypothesis tests for a mean. The difference is that this question has known \(\sigma\) and a one-sided alternative, so it uses a \(z\)-test instead of a \(t\)-test.

\section*{Mini-cheatsheet from this exam}

\subsection*{Probability}
Complement:
\[
P(A)=1-P(A^c).
\]
Conditional probability:
\[
P(A\mid B)=\frac{P(A\cap B)}{P(B)}.
\]
Union:
\[
P(A\cup B)=P(A)+P(B)-P(A\cap B).
\]

\subsection*{Continuous density}
To make a density:
\[
\int f(x)\,dx=1.
\]
CDF:
\[
F(x)=P(X\le x)=\int_{-\infty}^{x}f(t)\,dt.
\]
Right-tail probability:
\[
P(X>c)=1-F(c).
\]
Variance:
\[
V(X)=E(X^2)-[E(X)]^2.
\]

\subsection*{Poisson}
If \(X\sim \operatorname{Po}(\mu)\):
\[
P(X=k)=\frac{\mu^k}{k!}e^{-\mu}.
\]
\[
E(X)=\mu,\qquad V(X)=\mu.
\]
For independent variables:
\[
P(X_1\le c,\;X_2\le c)=P(X_1\le c)P(X_2\le c).
\]

\subsection*{Maximum likelihood}
Standard ML workflow:
\[
L(\theta)\rightarrow \ln L(\theta)\rightarrow \text{differentiate}\rightarrow \text{set equal to zero}.
\]

\subsection*{Variance and standard deviation}
Sample mean:
\[
\bar{x}=\frac{\sum x_i}{n}.
\]
Sample variance:
\[
s^2=\frac{\sum x_i^2-n\bar{x}^2}{n-1}.
\]
Standard deviation:
\[
s=\sqrt{s^2}.
\]

\subsection*{Confidence interval for \(\sigma\)}
For normal data:
\[
\left(
\sqrt{\frac{(n-1)s^2}{\chi^2_{0.95,n-1}}},
\sqrt{\frac{(n-1)s^2}{\chi^2_{0.05,n-1}}}
\right)
\]
for a 90\% interval.

\subsection*{Hypothesis test for mean, known \(\sigma\)}
\[
z=\frac{\bar{x}-\mu_0}{\sigma/\sqrt{n}}.
\]
Right-tailed 5\% test:
\[
\text{reject if } z>1.645.
\]

\subsection*{Power/sample size trick}
First write the rejection boundary, then calculate the probability of crossing that boundary under the new true mean.

\end{document}
